% Generated draft: requires mathematical and source review.
\documentclass[11pt]{article}
\usepackage[margin=25mm]{geometry}
\usepackage{fontspec}
\setmainfont{Latin Modern Roman}
\usepackage{amsmath,amssymb,mathtools,unicode-math}
\setmathfont{Latin Modern Math}
\usepackage{enumitem,needspace,xurl,hyperref}
\setlength{\parindent}{0pt}
\setlength{\parskip}{5pt}
\setlength{\emergencystretch}{3em}
\newcommand{\N}{\mathbb{N}}
\newcommand{\Z}{\mathbb{Z}}
\newcommand{\Q}{\mathbb{Q}}
\newcommand{\R}{\mathbb{R}}
\newcommand{\C}{\mathbb{C}}
\newcommand{\F}{\mathbb{F}}
\newcommand{\A}{\mathbb{A}}
\newcommand{\PP}{\mathbb{P}}
\newcommand{\E}{\mathbb{E}}
\newcommand{\eps}{\varepsilon}
\newcommand{\dd}{\,\mathrm{d}}
\begin{document}
\section*{UnsolvedMath ID: 800736}
\textbf{Source problem code:} catalogue entry
\subsection{Definitions and mathematical statement}
Let \(G\) be a topological group, with identity \(e\), and let \(\mathcal U_e(G)\) denote its open identity neighbourhoods. Write \(\N_0=\{0,1,2,\ldots\}\) and \(FU=\{fu:f\in F,\ u\in U\}\). Under the usual finite-translate convention, \(G\) is \(o\)-bounded if
\[
\forall (U_n)_{n\in\N_0}\in\mathcal U_e(G)^{\N_0}\quad
\exists (F_n)_{n\in\N_0},\quad
\bigl(\forall n\in\N_0,\ F_n\subseteq G\text{ is finite}\bigr)
\quad\text{and}\quad
G=\bigcup_{n\in\N_0}F_nU_n.
\]
The usual finite-translate interpretation of \(\mathrm{OF}(G)\) is the following perfect-information game of length \(\omega\). In inning \(n\), player I chooses \(U_n\in\mathcal U_e(G)\), and player II then chooses a finite set \(F_n\subseteq G\). Player II wins exactly when
\[
G=\bigcup_{n\in\N_0}F_nU_n;
\]
otherwise player I wins. A winning strategy for player I prescribes each neighbourhood from the preceding finite history and ensures failure of this covering equality against every legal sequence of responses by player II.

Under this interpretation, the question attributed to Banakh is: does there exist an \(o\)-bounded topological group \(G\) for which player I has a winning strategy in \(\mathrm{OF}(G)\)? Can such a group be chosen metrizable, meaning that its topology is induced by a metric?

\emph{Unresolved source qualifications.} The supplied record does not define \(o\)-boundedness or the rules of \(\mathrm{OF}(G)\). The definitions above give an explicit conventional interpretation, whose agreement with the cited source requires confirmation. The source writes ``(metrizable)'', leaving its intended relation to the unrestricted existence question implicit; both the unrestricted question and the metrizable strengthening are therefore displayed. The record also supplies no separation convention for topological groups. No Hausdorff or other additional separation hypothesis has been inserted; the intended convention must be checked before this formulation is accepted.
\subsection{Short English statement}
The problem asks whether an \(o\)-bounded topological group can nevertheless allow player I to force failure of the covering goal in \(\mathrm{OF}(G)\), and whether such an example can be metrizable. The finite-translate game interpretation and the intended separation assumptions require source confirmation.
\subsection{Sources}
\begin{itemize}
\item Local UnsolvedMath registry, permanent ID \(800736\), source problem code \texttt{catalogue entry}, in the supplied AMR Open Problem Lists record. The supplied status is \texttt{open}; this label is not an independently verified mathematical or literature assessment.
\item Banakh and Choban, \emph{Some Open Problems in Topological Algebra} (2012), Problem 10. Supplied source URL: \url{https://arxiv.org/abs/1202.1619}. The registry lists Banakh and Choban as proposers and gives the problem-specific attribution/context as Banakh.
\item The supplied provenance reports extraction by \texttt{arxiv\_latex} and an access date of 2026-10-16. This is supplied provenance, not a claim that the source was visited or independently checked in preparing this draft.
\end{itemize}
\end{document}
